%%
%% Template copyright (c) 2018-2024 Weitian LI
%% CC BY 4.0 License
%%

% 中文简历；请使用 XeLaTeX 编译。
\documentclass[zh]{resume}

% 最后一页页脚信息，可自行删除或修改。
\fileinfo{%
  \faEdit{} \today
}

% 中文模式下参数顺序为“名、姓”，最终显示为“况熙”。
\name{熙}{况}

\keywords{多模态算法, 大模型应用, AI 工程化, RAG, Agent, 计算机视觉}

\tagline{\texorpdfstring{\icon{\faBinoculars} }{}求职意向：多模态算法工程师 / 大模型应用开发工程师 / AI 工程化研发}

% 添加头像时取消下一行注释，并将 portrait.jpg 放在当前目录。
% \photo[shape=square, position=left]{6.5em}{portrait.jpg}

\profile{
  \mobile{199-8204-2910}
  \email{492700414@qq.com}
  \address{成都 / 北京均可}
  \\[0.2em]
  \icon{\faGraduationCap} 成都大学：计算机技术硕士（2024.09--至今） / 软件工程本科（2019.09--2023.06）
  \\[0.2em]
  \icon{\faAward} CET-6：468 \textbullet 成都大学学校奖学金一等奖 \textbullet 校学科建设交叉项目立项
  % 有 GitHub 主页后可取消注释：\github{your-account}
}

\begin{document}
\makeheader

%======================================================================
\begin{abstract}
计算机技术硕士在读，具备大模型应用、多模态检索、计算机视觉与后端工程化
落地经验。曾参与铁路配件智能识别、企业文档脱敏、医学报告质控等 AI 项目，
熟悉从模型适配、检索链路、业务接口到 Docker 部署的完整交付流程。技术栈覆盖
RAG、Agent、向量检索、图像分割、OCR、FastAPI、Milvus、MinIO、MySQL 等。
\end{abstract}

%======================================================================
\sectionTitle{技术能力}{\faWrench}
%======================================================================
\begin{competences}[9em]
  \comptence{大模型与智能应用}{%
    RAG、AI Agent、Prompt Engineering、LoRA 微调、GRPO 强化学习、
    LangChain、LlamaIndex
  }
  \comptence{多模态与视觉检索}{%
    SigLIP、DINOv2、BGE-M3、SAM/SAM2、YOLO-World、OCR、图像分割、
    图文检索、语义检索
  }
  \comptence{后端与工程化}{%
    Python、FastAPI、Flask、Django、MySQL、Redis、Milvus、MinIO、
    Docker、Linux、Nginx
  }
  \comptence{模型工程}{%
    PyTorch、模型推理部署、数据处理、向量索引构建、批处理任务、
    服务监控与日志排查
  }
\end{competences}

%======================================================================
\sectionTitle{实习经历}{\faBriefcase}
%======================================================================
\begin{experiences}
  \experience%
    [2026.03]%
    {2026.07}%
    {\textbf{布法罗机器人科技（成都）有限公司} \textbullet 多智能体研发实习生}%
    [\begin{itemize}
      \item 负责企业级 AI 多智能体架构搭建、流程编排与本地化部署落地
      \item 参与铁路配件智能识别系统核心研发，完成多模态检索、图像分割、
        向量索引、三维重构等业务能力建设
      \item 负责 SGO 业务系统、Kingbase/MySQL 数据源与 AI 服务接口适配，
        处理鉴权、数据同步、任务状态追踪和失败重试等交付问题
      \item 参与后端接口开发、业务联调、Docker 部署与甲方环境交付，
        支撑 AI 项目闭环上线
    \end{itemize}]

  \separator{0.8ex}
  \experience%
    [2025.11]%
    {2026.03}%
    {\textbf{四川安正企业管理集团} \textbullet 算法实习生}%
    [\begin{itemize}
      \item 负责法律文书 OCR、敏感信息抽取、敏感内容高亮与大模型 API 接入
      \item 参与 RAG 技术方案调研、知识库构建、文档分块、Embedding 选型
        与检索效果优化
      \item 完成后端业务接口封装、服务联调与调用监控建设，
        提升大模型应用可观测性
      \item 参与 Langfuse 调用链路接入与日志分析，支持大模型服务成本统计、
        异常定位和效果迭代
    \end{itemize}]
\end{experiences}

%======================================================================
\sectionTitle{科研经历}{\faFlask}
%======================================================================
\begin{experiences}
  \experience%
    [2024.06]%
    {至今}%
    {\textbf{医学影像报告质控与纠错方法研究}\\
     成都大学 \textbullet 硕士毕业论文}%
    [\begin{itemize}
      \item 面向医学影像报告质控场景，构建字符级表面错误纠错（CSC）
        与医学表达级语义纠错（CGEC）两类任务体系
      \item 构建 Medical-CSC-5500 数据集，覆盖少字、多字、错字和无错样本，
        基于 Qwen2.5-1.5B 完成医学领域继续预训练与 CSC LoRA 微调
      \item 引入加权 Levenshtein 距离与约束 Beam Search，融合语言模型概率、
        提示纠错概率和字符编辑代价，降低医学文本过纠与漏纠
      \item 在医学 CSC 测试集上，最终模型字符纠错 F1 达到 87.935\%，
        句子纠错 F1 达到 90.105\%，Segment F1 达到 92.523\%，FPR 降至 1.200\%
      \item 针对语义级表达错误，设计教师模型与 RAG 证据增强实验，
        比较无 RAG、RAG 和 Two-Stage RAG 的纠错效果
    \end{itemize}]
\end{experiences}

%======================================================================
\newpage
\sectionTitle{项目经历}{\faCode}
%======================================================================
\begingroup
\small
\setlength{\LTpre}{0pt}
\setlength{\LTpost}{0pt}
\begin{projects}
  \project%
    {2026.03}%
    [2026.07]%
    {布法罗机器人科技（成都）有限公司}%
    {工业多模态检索平台 \& AI 工程化交付}%
    {中车配件智能识别系统}%
    {%
      \begin{itemize}
        \item \textbf{系统研发：}参与从 0 到 1 搭建铁路配件智能识别平台，完成主业务服务、
          模型推理服务与三维重构服务解耦，支撑甲方内网部署、联调和验收
        \item \textbf{多模态检索：}融合 SigLIP、DINOv2、YOLO-World/SAM2、OCR 与 Milvus，
          打通图片预处理、特征提取、向量入库、相似召回、结果回写和业务 API 链路
        \item \textbf{询码语义检索：}基于 BGE-M3 构建历史询码文本索引，支持自然语言召回配件；
          设计全量重建、SGO 同步后增量索引、任务状态追踪与失败告警机制
        \item \textbf{企业系统对接：}对接 SGO 与 Kingbase/MySQL 数据源，处理配件主数据、图片、
          询码历史同步，以及私有时间类型兼容、失败明细、取消和重试等流程
        \item \textbf{工程交付：}完成 Docker 容器化交付，配置 MySQL、Redis、Milvus、MinIO、
          模型服务和重构服务依赖，实现 Token 鉴权、缓存和 Nginx + MinIO 文件代理访问
        \item \textbf{项目成果：}累计处理 7.8 万+ 配件主数据、1.4 万+ 样本图片、2 万+ 业务向量数据；
          建设 70+ 核心接口，覆盖配件管理、图文检索、技术核验、三维重构和文件代理闭环
      \end{itemize}
    }
  \separator{0.35ex}
  \project%
    {2025.11}%
    [2026.03]%
    {四川安正企业管理集团}%
    {大模型应用与文档智能}%
    {企业级文档脱敏一体机与大模型应用平台}%
    {%
      \begin{itemize}
        \item 基于 OCR、图像增强和大模型抽取姓名、公司、地址等 6 类敏感实体，
          完成敏感内容高亮、脱敏结果回写和业务接口封装
        \item 使用 LangChain 搭建文档分块、Embedding 向量化、知识检索与问答生成链路，
          完成 RAG Pipeline 整合与检索效果迭代
        \item 引入 Langfuse 搭建调用监控、日志追踪与用量统计体系，提升大模型服务排查、
          成本统计和运维效率
      \end{itemize}
    }
  \separator{0.35ex}
  \project%
    {2024.06}%
    [至今]%
    {成都大学}%
    {研究生课题}%
    {成都大学附属医院病历报告质控系统}%
    {%
      \begin{itemize}
        \item 面向医学报告实时与事后质控场景，构建影像--诊断问答数据集，
          基于 Qwen2.5-7B 完成 LoRA 微调并引入 GRPO 强化学习
        \item 构建医学知识图谱作为外部检索库，结合 RAG 提升质控结论的可解释性与准确性；
          使用 AI Agent 自动化构造质控规则，降低人工维护成本
        \item 相较基础微调模型，ROUGE-1 F1 提升 6.1\%，ROUGE-2 F1 提升 3.9\%
      \end{itemize}
    }
\end{projects}
\endgroup

\end{document}
